\documentclass[10pt,a4paper]{amsart}
\usepackage[margin=19mm]{geometry}
\usepackage{amsmath,amssymb,mathtools}
\usepackage[scr=rsfs,cal=euler]{mathalfa}
\usepackage{xfrac}
\usepackage[unicode,hidelinks,bookmarks=false]{hyperref}
\newcommand{\ie}{i.e.\ }
\newcommand{\cf}{cf.\ }
\newcommand{\resp}{respectively\ }
\newcommand{\Subsection}{\subsection}
\providecommand{\nobreakdash}{\nobreak\hbox{-}}
\setlength{\emergencystretch}{2em}
\multlinegap0pt
\usepackage{bm}
\usepackage{bbm}
\usepackage{dsfont}
\usepackage{xspace}
\usepackage{cancel}
\usepackage{mathtools}
\usepackage{amsthm}
\makeatletter
\@ifundefined{theorem}{\newtheorem{theorem}{Theorem}}{}
\makeatother
\newcommand{\abs}[1]{\left|{#1}\right|}
\newcommand{\bb}{\mathbb}
\newcommand{\bff}{\boldsymbol}
\newcommand{\ddt}{\frac{\mathrm{d}}{\mathrm{d}t}}
\newcommand{\ds}{\mathrm{d}s}
\newcommand{\inpro}[2]{\left\langle#1,#2\right\rangle}
\newcommand{\norm}[2]{\left\|{#1}\right\|_{#2}}
\title[Mixed finite element methods for the Landau--Lifshitz--Barya]{Mixed finite element methods for the Landau--Lifshitz--Baryakhtar and the regularised Landau--Lifshitz--Bloch equations in micromagnetics}
\author{Agus L. Soenjaya}
\date{Source: arXiv:2407.01125v1 (CC BY 4.0)}
\begin{document}
\maketitle

\noindent\emph{Source and licence.} Mixed finite element methods for the Landau--Lifshitz--Baryakhtar and the regularised Landau--Lifshitz--Bloch equations in micromagnetics. Version arXiv:2407.01125v1 (CC BY 4.0), distributed under the Creative Commons Attribution 4.0 International licence, as recorded in the arXiv licence metadata for this exact version and shown on the abstract page https://arxiv.org/abs/2407.01125v1. Full rights notice: licenses/arxiv-2407.01125-CC-BY-4.0.txt.\par\smallskip

\noindent\textit{Editorial scope.} Counted unit: the Theorem whose source label is \texttt{the:bloch to bar} in the article (source0.tex lines 6498--6509, source label \texttt{the:bloch to bar}), delivered together with the article's own complete original proof of it (source0.tex lines 6511--6830, one \texttt{proof} environment of 320 lines). No mathematics is written, repaired or re-derived: statement and proof are the article's own text line for line. Required context retained uncounted: equ:u H bloch H2 (lines 6422--6425); equ:ue est (lines 6429--6439). Every definition, display and label of those lines is retained unchanged. Omitted from this package and not redistributed: 1--6421, 6426--6428, 6440--6497, 6510--6510, 6831--7514. Nothing inside a retained excerpt is omitted or altered: every retained body is delivered exactly as the article's lines are written, so no omission receipt of any kind accompanies this package. Citations: the retained excerpts carry no citation command at all, so no bibliography entry of the article is transcribed and no reference is redistributed. Reference adaptation: none was needed and none was made; every reference inside the delivered unit resolves inside it, so no unresolved reference or citation is produced. Printed slips kept literal: no printed word, symbol, hypothesis or number of the counted statement or of its proof was altered. Layout and class: the article's own class is \texttt{amsart} with packages including amsmath, amsthm, amscd, amsfonts, amssymb, graphicx, color, mathtools, mathrsfs, hyperref, enumerate, setspace, multicol, geometry; the wrapper is the lane's pinned \texttt{amsart} editorial wrapper at 10pt on A4, which restates the theorem-like environments the retained text needs and adds the article's own macro definitions that the retained text needs (\texttt{\textbackslash abs}, \texttt{\textbackslash bb}, \texttt{\textbackslash bff}, \texttt{\textbackslash ddt}, \texttt{\textbackslash ds}, \texttt{\textbackslash inpro}, \texttt{\textbackslash norm}), with extra wrapper packages (bm, bbm, dsfont, xspace, cancel, mathtools). The delivered numbering is the unit's own. Assets: no retained span contains a figure environment, a graphics-inclusion command, a PDF-page inclusion, a table or a TikZ picture; no image file is copied into this package, the package declares no asset dependency and no image file is redistributed with it. Licence: version arXiv:2407.01125v1 of this article is distributed under the Creative Commons Attribution 4.0 International licence, as recorded in the arXiv licence metadata for this exact version and shown on the abstract page https://arxiv.org/abs/2407.01125v1. A full rights notice is delivered at licenses/arxiv-2407.01125-CC-BY-4.0.txt. Characters: every retained excerpt is pure ASCII, so the delivered engine, XeLaTeX with its default Latin Modern text font, reports no missing character.
\begin{align}\label{equ:u H bloch H2}
	\norm{\bff{u}}{L^\infty(\bb{H}^2)} + \norm{\bff{u}}{L^2(\bb{H}^3)} +
	\norm{\partial_t \bff{u}}{L^\infty(\bb{L}^2)} + \norm{\bff{H}}{L^2(\bb{H}^1)} \leq K_2,
\end{align}

\begin{align}\label{equ:ue est}
	\norm{\bff{u}^\varepsilon}{L^\infty(\bb{H}^1)} 
	+
	\norm{\Delta \bff{u}^\varepsilon}{L^2(\bb{L}^2)}
	+
	\norm{\bff{H}^\varepsilon}{L^2(\bb{L}^2)}
	+
	\sqrt{\varepsilon} \norm{\nabla \bff{H}^\varepsilon}{L^2(\bb{L}^2)}
	\leq
	K_3.
\end{align}

\begin{theorem}\label{the:bloch to bar}
Let $\bff{u}^\varepsilon$ be a strong solution of the LLBar equation with $\lambda_e=\varepsilon$ and initial data $\bff{u}_0^\varepsilon\in \bb{H}^2$. Let $\bff{u}$ be a strong solution of the LLBloch equation ($\lambda_e=0$) with initial data~$\bff{u}_0\in \bb{H}^2$. Then
\begin{align}\label{equ:est ue u H1}
	\norm{\bff{u}^\varepsilon-\bff{u}}{L^\infty(\bb{H}^1)}
	+
	\norm{\bff{u}^\varepsilon-\bff{u}}{L^2(\bb{H}^2)}
	+
	\norm{\bff{H}^\varepsilon-\bff{H}}{L^2(\bb{L}^2)}
	&\leq
	C \sqrt{\varepsilon}.
\end{align}
\end{theorem}

\begin{proof}
We write $\bff{v}:= \bff{u}^\varepsilon-\bff{u}$ and $\bff{B}:= \bff{H}^\varepsilon-\bff{H}$. Let $\bff{v}_0:= \bff{u}_0^\varepsilon-\bff{u}_0$, so that $\bff{v}(0)=\bff{v}_0$. Since $\bff{u}^\varepsilon$ and $\bff{u}$ are strong solutions of the corresponding equations, we have
\begin{equation}\label{equ:diff v}
	\begin{aligned}
		\partial_t \bff{v}
		&=
		\lambda_r \bff{B}
		- \varepsilon \Delta \bff{H}^\varepsilon 
		- \gamma \bff{u}^\varepsilon \times \bff{B}
		- \gamma \bff{v}\times \bff{H},
		\\
		\bff{B}
		&=
		\Delta \bff{v}
		+ \kappa\mu \bff{v}
		-
		\kappa |\bff{u}^\varepsilon|^2 \bff{v}
		-
		\kappa \big( (\bff{u}^\varepsilon+\bff{u})\cdot \bff{v} \big) \bff{u}
		-
		\beta \bff{e}(\bff{e}\cdot \bff{v}).
	\end{aligned}
\end{equation}
%Taking inner product of both equations in~\eqref{equ:diff v} with $\bff{v}$, we have
%\begin{equation}\label{equ:diff v dot v}
%	\begin{aligned}
%		\frac{1}{2} \ddt \norm{\bff{v}}{\bb{L}^2}^2
%		&=
%		\lambda_r \inpro{\bff{B}}{\bff{v}} 
%		+
%		\varepsilon \inpro{\nabla \bff{H}^\varepsilon}{\nabla \bff{v}} 
%		- 
%		\gamma \inpro{\bff{u}^\varepsilon \times \bff{B}}{\bff{v}},
%		\\
%		\inpro{\bff{B}}{\bff{v}}
%		&=
%		-
%		\norm{\nabla \bff{v}}{\bb{L}^2}^2 
%		+ 
%		\kappa\mu \norm{\bff{v}}{\bb{L}^2}^2 
%		-
%		\kappa \norm{|\bff{u}^\varepsilon|\, |\bff{v}|}{\bb{L}^2}^2 
%		-
%		\kappa \inpro{(\bff{u}^\varepsilon \cdot \bff{v})\bff{u}}{\bff{v}}
%		-
%		\kappa \norm{\bff{u}\cdot \bff{v}}{L^2}^2
%		-
%		\beta \norm{\bff{e}\cdot \bff{v}}{L^2}^2.
%	\end{aligned}
%\end{equation}
%We note that
%\begin{equation}\label{equ:inpro ue B v}
%	\gamma\inpro{\bff{u}^\varepsilon \times \bff{B}}{\bff{v}}
%	=
%	-\gamma \inpro{\bff{v}\times \nabla \bff{u}^\varepsilon}{\nabla \bff{v}}
%	-\kappa\gamma \inpro{\bff{u}^\varepsilon \times \big((\bff{u}^\varepsilon+\bff{u})\cdot \bff{v}\big)\bff{u}}{\bff{v}}
%	-
%	\beta\gamma \inpro{\bff{u}^\varepsilon \times \bff{e}(\bff{e}\cdot \bff{v})}{\bff{v}}.
%\end{equation}
%Substituting \eqref{equ:inpro ue B v} into~\eqref{equ:diff v dot v}, we obtain
%\begin{align}\label{equ:v L2}
%	&\frac{1}{2} \ddt \norm{\bff{v}}{\bb{L}^2}^2
%	+
%	\lambda_r \norm{\nabla \bff{v}}{\bb{L}^2}^2
%	+
%	\kappa \norm{|\bff{u}^\varepsilon|\, |\bff{v}|}{\bb{L}^2}^2
%	+
%	\kappa \norm{\bff{u}\cdot \bff{v}}{L^2}^2
%	+
%	\beta \norm{\bff{e}\cdot \bff{v}}{L^2}^2 
%	-
%	\kappa\mu \norm{\bff{v}}{\bb{L}^2}^2
%	\nonumber \\
%	&=
%	-
%	\kappa \inpro{(\bff{u}^\varepsilon\cdot \bff{v})\bff{u}}{\bff{v}}
%	-
%	\varepsilon \inpro{\nabla \bff{H}^\varepsilon}{\nabla \bff{v}}
%	+
%	\gamma \inpro{\bff{v}\times \nabla \bff{u}^\varepsilon}{\nabla \bff{v}}
%	+
%	\kappa\gamma \inpro{\bff{u}^\varepsilon \times \big((\bff{u}^\varepsilon+\bff{u})\cdot \bff{v}\big)\bff{u}}{\bff{v}}
%	\nonumber \\
%	&\quad
%	+
%	\beta\gamma \inpro{\bff{u}^\varepsilon \times \bff{e}(\bff{e}\cdot \bff{v})}{\bff{v}}
%	\nonumber \\
%	&=:
%	I_1 + I_2+ I_3 + I_4+I_5.
%\end{align}
%It remains to estimate each term on the last line, noting estimates~\eqref{equ:u H bloch H1} and~\eqref{equ:ue est}. Firstly, by Holder's and Young's inequalities and the Sobolev embedding $\bb{H}^1\hookrightarrow \bb{L}^6$,
%\begin{align*}
%	\abs{I_1}
%	\leq
%	\kappa \norm{\bff{u}^\varepsilon}{\bb{L}^6}
%	\norm{\bff{v}}{\bb{L}^6}
%	\norm{\bff{u}}{\bb{L}^6}
%	\norm{\bff{v}}{\bb{L}^2}
%	\leq
%	C\norm{\bff{v}}{\bb{L}^2}^2
%	+
%	\frac{\lambda_r}{8} \norm{\nabla \bff{v}}{\bb{L}^2}^2.
%\end{align*}
%For the term $I_2$, by Young's inequality,
%\begin{align*}
%	\abs{I_2}
%	\leq
%	\frac{2\varepsilon^2}{\lambda_r} \norm{\nabla \bff{H}^\varepsilon}{\bb{L}^2}^2
%	+
%	\frac{\lambda_r}{8} \norm{\nabla \bff{v}}{\bb{L}^2}^2.
%\end{align*}
%For the terms $I_3$, $I_4$, and $I_5$, by similar argument we have
%\begin{align*}
%	\abs{I_3}
%	&\leq
%	\gamma \norm{\bff{v}}{\bb{L}^2} \norm{\nabla \bff{u}^\varepsilon}{\bb{L}^\infty} \norm{\nabla \bff{v}}{\bb{L}^2}
%	\leq
%	C \norm{\nabla \bff{u}^\varepsilon}{\bb{L}^\infty}^2 \norm{\bff{v}}{\bb{L}^2}^2
%	+
%	\frac{\lambda_r}{8} \norm{\nabla \bff{v}}{\bb{L}^2}^2
%	\\
%	\abs{I_4}
%	&\leq
%	\norm{\bff{u}^\varepsilon}{\bb{L}^6}
%	\norm{\bff{u}^\varepsilon+\bff{u}}{\bb{L}^6}
%	\norm{\bff{v}}{\bb{L}^6}
%	\norm{\bff{u}}{\bb{L}^\infty}
%	\norm{\bff{v}}{\bb{L}^2}
%	\leq
%	C\norm{\bff{v}}{\bb{L}^2}^2
%	+
%	\frac{\lambda_r}{8} \norm{\nabla \bff{v}}{\bb{L}^2}^2,
%	\\
%	\abs{I_5}
%	&\leq
%	\beta\gamma \norm{\bff{u}^\varepsilon}{\bb{L}^3}
%	\norm{\bff{v}}{\bb{L}^6} \norm{\bff{v}}{\bb{L}^2}
%	\leq
%	C\norm{\bff{v}}{\bb{L}^2}^2
%	+
%	\frac{\lambda_r}{8} \norm{\nabla \bff{v}}{\bb{L}^2}^2,
%\end{align*}
%where in the above, the constant $C$ is independent of $\varepsilon$. 
%Altogether, substituting these estimates into~\eqref{equ:v L2} and integrating, we obtain
%\begin{align*}
%	\norm{\bff{v}(t)}{\bb{L}^2}^2 
%	+
%	\int_0^t \norm{\nabla \bff{v}(s)}{\bb{L}^2}^2\,\ds
%	\leq
%	C \int_0^t \norm{\bff{v}(s)}{\bb{L}^2}^2 \,\ds
%	+
%	\varepsilon^2 \int_0^t \norm{\nabla \bff{H}^\varepsilon}{\bb{L}^2}^2
%	\leq
%	C \int_0^t \norm{\bff{v}(s)}{\bb{L}^2}^2 \,\ds
%	+
%	C_3 \varepsilon,
%\end{align*}
%where $C_3$ is the constant in~\eqref{equ:ue est}. Invoking the Gronwall inequality yields~\eqref{equ:est ue u L2}. Uniqueness of the solution to the LLBloch equation follows from the corresponding uniqueness result for the LLBar equation~\cite{SoeTra23}.
Successively taking the inner product of the first equation in~\eqref{equ:diff v} with $\bff{B}$, then taking the inner product of the second equation in~\eqref{equ:diff v} with $-\partial_t \bff{v}$, we obtain
\begin{align*}
	\inpro{\partial_t \bff{v}}{\bff{B}}
	&=
	\lambda_r \norm{\bff{B}}{\bb{L}^2}^2
	+
	\varepsilon \norm{\nabla \bff{H}^\varepsilon}{\bb{L}^2}^2
	-
	\varepsilon \inpro{\nabla \bff{H}^\varepsilon}{\nabla \bff{H}}
	-
	\gamma \inpro{\bff{v}\times \bff{H}}{\bff{B}}
	\\
	-\inpro{\bff{B}}{\partial_t \bff{v}}
	&=
	\frac12 \ddt \norm{\nabla \bff{v}}{\bb{L}^2}^2
	-
	\frac{\kappa \mu}{2} \ddt \norm{\bff{v}}{\bb{L}^2}^2
	+
	\kappa \inpro{|\bff{u}^\varepsilon|^2 \bff{v}}{\partial_t \bff{v}}
	+
	\kappa \inpro{\big( (\bff{u}^\varepsilon+\bff{u})\cdot \bff{v} \big) \bff{u}}{\partial_t \bff{v}}
	+
	\frac{\beta}{2} \norm{\bff{e}\cdot \bff{v}}{L^2}^2,
\end{align*}
where on the first equation we integrated by parts and used $\inpro{\nabla \bff{H}^\varepsilon}{\nabla \bff{B}}= \norm{\nabla \bff{H}^\varepsilon}{\bb{L}^2}^2 - \inpro{\nabla \bff{H}^\varepsilon}{\nabla \bff{H}}$.
Adding the above equations (noting $\mu<0$) gives
\begin{align}\label{equ:add dtv B}
	&\frac12 \ddt \norm{\nabla \bff{v}}{\bb{L}^2}^2
	-
	\frac{\kappa\mu}{2} \ddt \norm{\bff{v}}{\bb{L}^2}^2
	+
	\frac{\beta}{2} \norm{\bff{e}\cdot \bff{v}}{L^2}^2
	+
	\lambda_r \norm{\bff{B}}{\bb{L}^2}^2
	+
	\varepsilon \norm{\nabla \bff{H}^\varepsilon}{\bb{L}^2}^2
	\nonumber \\
	&=
	\varepsilon \inpro{\nabla \bff{H}^\varepsilon}{\nabla \bff{H}}
	+
	\gamma \inpro{\bff{v}\times \bff{H}}{\bff{B}}
	-
	\kappa \inpro{|\bff{u}^\varepsilon|^2 \bff{v}}{\partial_t \bff{v}}
	-
	\kappa \inpro{\big( (\bff{u}^\varepsilon+\bff{u})\cdot \bff{v} \big) \bff{u}}{\partial_t \bff{v}}
	\nonumber \\
	&=: J_1+J_2+J_3+J_4.
\end{align}
We will estimate each term on the last line. In the following, the constant $C$ is independent of $\varepsilon$. Firstly, by Young's inequality,
\begin{align}\label{equ:J1}
	\abs{J_1}
	&\leq
	\frac{\varepsilon}{4} \norm{\nabla \bff{H}^\varepsilon}{\bb{L}^2}^2
	+
	4\varepsilon \norm{\nabla \bff{H}}{\bb{L}^2}^2.
\end{align}
Next, by Young's inequality and Sobolev embedding,
\begin{align}\label{equ:J2}
	\abs{J_2}
	&\leq
	\frac{\lambda_r}{8} \norm{\bff{B}}{\bb{L}^2}^2
	+
	C \norm{\bff{H}}{\bb{L}^4}^2 \norm{\bff{v}}{\bb{L}^4}^2
	\leq
	\frac{\lambda_r}{8} \norm{\bff{B}}{\bb{L}^2}^2
	+
	C \norm{\bff{H}}{\bb{H}^1}^2 \norm{\bff{v}}{\bb{H}^1}^2.
\end{align}
We now aim to estimate $J_3$. To this end, substituting $\partial_t \bff{v}$ by the first equation in~\eqref{equ:diff v} and integrating by parts as necessary, we have
\begin{align}\label{equ:L123}
	 J_3
	 =
	 -\kappa\lambda_r \inpro{|\bff{u}^\varepsilon|^2 \bff{v}}{\bff{B}}
	 +
	 \kappa \varepsilon \inpro{\nabla\big(|\bff{u}^\varepsilon|^2 \bff{v}\big)}{\nabla \bff{H}^\varepsilon}
	 +
	 \kappa \gamma \inpro{|\bff{u}^\varepsilon|^2 \bff{v}}{\bff{u}^\varepsilon\times \bff{B}}
	 =: J_{3a}+J_{3b}+J_{3c}.
\end{align}
For the terms $J_{3a}$ and $J_{3b}$, by Young's inequality and the Sobolev embedding (noting~\eqref{equ:ue est}) we have
\begin{align*}
	\abs{J_{3a}} 
	\leq
	\kappa \lambda_r \norm{\bff{u}^\varepsilon}{\bb{L}^6}^2 \norm{\bff{v}}{\bb{L}^6} \norm{\bff{B}}{\bb{L}^2}
	\leq
	\frac{\lambda_r}{8} \norm{\bff{B}}{\bb{L}^2}
	+
	C \norm{\bff{v}}{\bb{H}^1}^2.
\end{align*}
For the terms $J_{3b}$ and $J_{3c}$, similarly we obtain
\begin{align*}
	\abs{J_{3b}}
	&\leq
	\frac{\varepsilon}{4} \norm{\nabla \bff{H}^\varepsilon}{\bb{L}^2}^2
	+
	4\varepsilon \norm{\bff{u}^\varepsilon}{\bb{L}^6}^2 \norm{\nabla \bff{u}^\varepsilon}{\bb{L}^6}^2 \norm{\bff{v}}{\bb{L}^6}^2
	+
	4 \varepsilon \norm{\bff{u}^\varepsilon}{\bb{L}^6}^4 \norm{\nabla \bff{v}}{\bb{L}^6}^2
	\\
	&\leq
	\frac{\varepsilon}{4} \norm{\nabla \bff{H}^\varepsilon}{\bb{L}^2}^2
	+
	C\varepsilon \norm{\Delta \bff{u}^\varepsilon}{\bb{L}^2}^2 \norm{\bff{v}}{\bb{H}^1}^2
	+
	C\varepsilon \norm{\Delta \bff{v}}{\bb{L}^2}^2,
	\\
	\abs{J_{3c}}
	&\leq
	\kappa \gamma \norm{\bff{u}^\varepsilon}{\bb{L}^6}^2 \norm{\bff{v}}{\bb{L}^6}
	\norm{\bff{u}^\varepsilon}{\bb{L}^\infty} \norm{\bff{B}}{\bb{L}^2}
	\leq
	\frac{\lambda_r}{8} \norm{\bff{B}}{\bb{L}^2}
	+
	\norm{\bff{u}^\varepsilon}{\bb{H}^2}^2 \norm{\bff{v}}{\bb{H}^1}^2.
\end{align*}
Substituting these estimates into~\eqref{equ:L123}, we obtain
\begin{align}\label{equ:J3}
	\abs{J_3}
	\leq
	\frac{\lambda_r}{4} \norm{\bff{B}}{\bb{L}^2}^2
	+
	\frac{\varepsilon}{4} \norm{\nabla \bff{H}^\varepsilon}{\bb{L}^2}^2
	+
	C\varepsilon \norm{\Delta \bff{v}}{\bb{L}^2}^2
	+
	C\left(1+\norm{\bff{u}^\varepsilon}{\bb{H}^2}^2 \right)\norm{\bff{v}}{\bb{H}^1}^2,
\end{align}
where $C$ is independent of $\varepsilon$.
The term $J_4$ can be estimated in a similar manner as $J_3$, resulting in
\begin{align}\label{equ:J4}
	\abs{J_4}
	\leq
	\frac{\lambda_r}{4} \norm{\bff{B}}{\bb{L}^2}^2
	+
	\frac{\varepsilon}{4} \norm{\nabla \bff{H}^\varepsilon}{\bb{L}^2}^2
	+
	C\varepsilon \norm{\Delta \bff{v}}{\bb{L}^2}^2
	+
	C\left(1+\norm{\bff{u}^\varepsilon}{\bb{H}^2}^2 + \norm{\bff{u}}{\bb{H}^2}^2
	+ \norm{\bff{H}}{\bb{L}^2}^2 \right)\norm{\bff{v}}{\bb{H}^1}^2.
\end{align}
Altogether, substituting the estimates~\eqref{equ:J1}, \eqref{equ:J2}, \eqref{equ:J3}, and~\eqref{equ:J4} into~\eqref{equ:add dtv B}, taking care to absorb relevant terms to the left-hand side, we obtain
\begin{align*}
	\ddt \norm{\bff{v}}{\bb{H}^1}^2
	+
	\norm{\bff{B}}{\bb{L}^2}^2
	\leq
	C\varepsilon \norm{\nabla \bff{H}}{\bb{L}^2}^2
	+
	C\varepsilon \norm{\Delta \bff{v}}{\bb{L}^2}^2
	+
	C\left(1+\norm{\bff{u}^\varepsilon}{\bb{H}^2}^2 + \norm{\bff{u}}{\bb{H}^2}^2
	+ \norm{\bff{H}}{\bb{H}^1}^2 \right)\norm{\bff{v}}{\bb{H}^1}^2.
\end{align*}
We now integrate both sides with respect to $t$. Note that by~\eqref{equ:u H bloch H2} and~\eqref{equ:ue est}, we have
\begin{align*}
	\int_0^t \left(1+ \norm{\bff{u}^\varepsilon(s)}{\bb{H}^2}^2 + \norm{\bff{u}(s)}{\bb{H}^2}^2
	+ \norm{\bff{H}(s)}{\bb{H}^1}^2 \right) \ds \leq 
	T+K_2+K_3.
\end{align*}
Invoking the Gronwall inequality (noting~\eqref{equ:u H bloch H2} and~\eqref{equ:ue est} again), we obtain~\eqref{equ:est ue u H1}. This completes the proof of the theorem.
\end{proof}

\end{document}
